\lab{10}{Цитатный поиск и координатный индекс}
\task
Расширить булев поиск точными цитатами и упорядоченными цитатами
\code{"..." / N}, где расстояние между первым и последним словом
не превышает \(N\). Описать координатный формат, размер, число
позиций, средние количества позиций, скорость и примеры запросов.

\method
Для каждой пары «терм, документ» хранится возрастающий список позиций
вхождений. Позиция - номер токена после токенизации, а не смещение
символа. Стоп-слова не удаляются, поэтому координаты соседних слов
сохраняют исходное расстояние. Нормализация меняет текст терма,
но не количество его вхождений.

Кандидаты цитаты должны содержать все её термы. В качестве ведущего
берётся терм с минимальным \(df\); курсоры остальных списков продвигаются
к текущему документу. Для точной цитаты длины \(m\) проверяются
позиции \(p,p+1,\ldots,p+m-1\). Для цитаты с расстоянием
выбираются строго возрастающие позиции \(p_1<\cdots<p_m\),
удовлетворяющие \(p_m-p_1\le N\). Повтор слова в цитате требует
разных последовательных вхождений. Порядок слов не переставляется.

Текущий формат - \code{IRPOS003}. Общий заголовок занимает 144 байта.
Первые 96 байт имеют назначение, аналогичное булевому формату,
но сигнатура и версия другие. Продолжение:
\begin{table}[H]\centering
\begin{tabular}{rll}\toprule
Смещение & Размер & Поле\\\midrule
96 & 8 & Секция длин документов\\
104 & 8 & Секция прыжков\\
112 & 8 & Секция норм векторов\\
120 & 8 & Отпечаток корпуса\\
128 & 4 & Шаг прыжков\\
132 & 4 & Идентификатор словаря WordNet\\
136 & 8 & Размер заголовка: 144\\\bottomrule
\end{tabular}
\caption{Расширение заголовка координатного индекса.}
\end{table}
Терм занимает 56 байт: пять 8-байтовых полей (смещение текста,
смещение постингов, длина постингов, смещение прыжков, суммарная
частота) и четыре 4-байтовых поля (длина текста, \(df\), число
блоков, резерв). Постинг включает \(docID\), число вхождений и
позиции. В raw-режиме каждое число имеет размер 4 байта;
VB-режим описан в следующей работе.

Прямой индекс дополнен длинами документов и нормами для TF-IDF.
Отпечаток зависит от URL, заголовков и текстов; он позволяет
проверять согласованность индексов одного корпуса. Это служебный
некриптографический отпечаток, не доказательство отсутствия коллизий.

\journal
Пары булевой индексации расширены координатами. Добавлены внешняя
сортировка позиций, агрегирование по терму и документу, курсор
постинг-листа и проверка цитат. Разбор кавычек и \code{/ N}
встроен в общий парсер. Дистанция сравнивается в расширенном
целом типе, чтобы сумма позиции и длины фразы не переполнялась.

\testing
Для документов \code{alpha beta beta gamma}, \code{alpha x beta gamma},
\code{beta alpha gamma} проверяются точный порядок и разрешённые
вкрапления. Запрос \code{"alpha beta"} не должен находить второй
документ, а \code{"alpha beta" / 2} должен. Повторная цитата
\code{"alpha alpha"} проверяет разные вхождения. Незакрытая
кавычка и отрицательная дистанция должны вызывать ошибку.
Все четыре теста координатного набора повторно пройдены.

\results
Для полного корпуса получены 20\,430\,767 термов,
640\,214\,849 пар «терм, документ» и 2\,483\,697\,508 позиций.
Среднее число позиций на терм - 121,5665, на постинг - 3,8795.
Вариант raw без прыжков занимает 16\,694\,523\,464 байта,
координатная часть - 15\,056\,508\,824 байта. Его
сборка заняла 2214,25 секунды (36,90 минуты), создала 29 блоков,
скорость составила 7601,84 KiB/с. Пиковый RSS составил
1\,230\,024 KiB, около 1201,20 MiB. Эти времена относятся именно
к raw-сборке из эксперимента №11.

\begin{table}[H]\centering
\begin{tabular}{lrr}\toprule
Запрос & Совпадений & Время, мс\\\midrule
\code{memory allocation} & 36\,917 & 233,170\\
\code{mutex locking} & 3\,230 & 24,931\\
\code{workqueues} & 43 & 0,106\\
\code{mprotect} & 621 & 0,325\\
\code{file descriptors} & 16\,623 & 234,695\\
\code{kernel sanitizers} & 35 & 3,519\\
\code{epolll} & 0 & 0,067\\
\bottomrule\end{tabular}
\caption{Одиночные поисковые замеры координатного движка по полному корпусу; это не среднее нескольких повторов.}
\end{table}

В таблице приведены обычные строгие запросы координатного движка.
Запрос \code{"memory allocation"} на полном корпусе находит
16\,449 документов, а \code{"memory allocation" / 5} находит
17\,815 документов. Расширение дистанции допускает вкрапления
между словами и увеличивает множество совпадений.
Проверены совпадение выдачи raw- и VB-индексов, вариантов с разным
шагом прыжков, а также поиск без доступа к исходному корпусу.
Цитата может сочетаться с фильтром, например
\code{"memory allocation" && !"out of memory"}.

\conclusions
Координаты позволяют отличить совместное присутствие слов от фразы.
Цена - значительно больший объём: сохраняются все вхождения.
Частые слова и длинные списки позиций создают дорогие запросы.
Для ускорения полезны ведущий редкий терм, отбрасывание неподходящих
документов до проверки координат, сжатие и прыжки. «Точная цитата»
здесь точна по нормализованным токенам, а не по исходной строке:
регистр, пунктуация и словоформы могут различаться.
